\chapter{Particle-in-a-Box Residuals and Identifiability}
\index{particle in a box|textbf}
\index{identifiability!particle in a box}

\label{app:particle-box}

\section{Purpose of the example}

The one-dimensional particle in a box provides a minimal setting in which the
continuum Hamiltonian, boundary conditions, discrete representation, spectral
restriction, and effective-model interpretation can be separated explicitly.
The example is useful not because the spectrum is difficult to obtain, but
because every state-space map can be written in closed form. It therefore
exposes several distinct residual operators that can otherwise be conflated
under the notation $V_{\mathrm{eff}}$.

\section{Continuum Hamiltonian}

Let a particle of mass $m$ be confined to the interval
\begin{equation}
    \Omega=(0,L).
\end{equation}
The Hamiltonian is the Dirichlet realization of the kinetic-energy differential
expression,
\begin{equation}
    \hat H_{\mathrm{box}}
    =
    -\frac{\hbar^2}{2m}
    \frac{d^2}{dx^2},
\end{equation}
with domain
\begin{equation}
    \mathcal D(\hat H_{\mathrm{box}})
    =
    H^2(0,L)\cap H_0^1(0,L).
\end{equation}
Equivalently,
\begin{equation}
    \psi(0)=\psi(L)=0.
\end{equation}
The eigenfunctions and eigenvalues are
\begin{equation}
    \psi_n(x)
    =
    \sqrt{\frac{2}{L}}
    \sin\left(\frac{n\pi x}{L}\right),
\end{equation}
and
\begin{equation}
    E_n
    =
    \frac{\hbar^2\pi^2n^2}{2mL^2},
    \qquad
    n=1,2,\ldots.
\end{equation}
The interior potential is zero. Confinement is encoded in the domain of the
self-adjoint kinetic operator rather than by finite values of a multiplication
potential inside the interval~\cite{bonneau2001,angelone2024}. A different
boundary condition applied to the same formal differential expression defines
a different operator.

\section{Finite-difference realization}
\index{finite difference}


Introduce $N$ interior points
\begin{equation}
    x_j=j\Delta x,
    \qquad
    \Delta x=\frac{L}{N+1},
    \qquad
    j=1,\ldots,N.
\end{equation}
The boundary values are fixed to zero and are not independent components of
the numerical state vector. With the centered second-order approximation,
\begin{equation}
    \left.
        \frac{d^2\psi}{dx^2}
    \right|_{x=x_j}
    \approx
    \frac{
        \psi(x_{j+1})-2\psi(x_j)+\psi(x_{j-1})
    }{\Delta x^2},
\end{equation}
the discrete Hamiltonian is
\begin{equation}
    \mathbf H_h
    =
    \frac{\hbar^2}{2m\Delta x^2}
    \begin{bmatrix}
        2 & -1 & 0 & \cdots & 0 \\
        -1 & 2 & -1 & \ddots & \vdots \\
        0 & -1 & 2 & \ddots & 0 \\
        \vdots & \ddots & \ddots & \ddots & -1 \\
        0 & \cdots & 0 & -1 & 2
    \end{bmatrix}.
\end{equation}
This matrix is the second-order discrete Dirichlet Laplacian multiplied by the
kinetic prefactor~\cite{strikwerda2004}. Its eigenvalues are
\begin{equation}
    \varepsilon_i
    =
    \frac{2\hbar^2}{m\Delta x^2}
    \sin^2\left(
        \frac{i\pi}{2(N+1)}
    \right),
    \qquad
    i=1,\ldots,N.
\end{equation}
For fixed $i$,
\begin{equation}
    \varepsilon_i\longrightarrow E_i
\end{equation}
as $N\rightarrow\infty$. Writing
\begin{equation}
    z_i=\frac{i\pi}{2(N+1)},
\end{equation}
the dispersion error satisfies
\begin{equation}
    \frac{\varepsilon_i}{E_i}
    =
    \left(\frac{\sin z_i}{z_i}\right)^2,
\end{equation}
and hence
\begin{equation}
    \frac{\varepsilon_i-E_i}{E_i}
    =
    -\frac{z_i^2}{3}+O(z_i^4).
\end{equation}
The relation between a Euclidean-normalized eigenvector $\mathbf v_i$ and
samples of a normalized continuum eigenfunction contains the quadrature factor
\begin{equation}
    v_i(j)
    \approx
    \sqrt{\Delta x}\,\psi_i(x_j),
\end{equation}
up to an arbitrary phase. This factor distinguishes the numerical inner
product from the continuum $L^2$ inner product.

\section{Spectral reconstruction and retention}

Let
\begin{equation}
    \boldsymbol\Phi
    =
    \begin{bmatrix}
        \mathbf v_1 & \mathbf v_2 & \cdots & \mathbf v_N
    \end{bmatrix},
\end{equation}
and
\begin{equation}
    \boldsymbol\Lambda
    =
    \operatorname{diag}
    (\varepsilon_1,\ldots,\varepsilon_N).
\end{equation}
The discrete Hamiltonian has the exact spectral reconstruction
\begin{equation}
    \mathbf H_h
    =
    \boldsymbol\Phi
    \boldsymbol\Lambda
    \boldsymbol\Phi^\dagger.
\end{equation}
This is an exact identity for the discrete matrix. It is not an equality
between the finite-difference matrix and the continuum differential operator.

Retain the lowest $M$ discrete eigenvectors in
\begin{equation}
    \boldsymbol\Phi_M
    =
    \begin{bmatrix}
        \mathbf v_1 & \cdots & \mathbf v_M
    \end{bmatrix},
\end{equation}
and define
\begin{equation}
    \mathbf P_M
    =
    \boldsymbol\Phi_M
    \boldsymbol\Phi_M^\dagger.
\end{equation}
The retained Hamiltonian embedded in the full grid space is
\begin{equation}
    \mathbf H_h^{(P_M)}
    =
    \mathbf P_M\mathbf H_h\mathbf P_M
    =
    \boldsymbol\Phi_M
    \boldsymbol\Lambda_M
    \boldsymbol\Phi_M^\dagger,
\end{equation}
whereas its representation in retained coordinates is
\begin{equation}
    \mathbf H_{\mathrm{red}}
    =
    \boldsymbol\Phi_M^\dagger
    \mathbf H_h
    \boldsymbol\Phi_M
    =
    \boldsymbol\Lambda_M.
\end{equation}
The first matrix is $N\times N$ and has rank $M$; the second is $M\times M$.
They represent the same retained operator in different coordinate spaces.
Their dimensions and embeddings must be accounted for before subtraction.

\section{Three inequivalent residuals}
\index{residual!inequivalent constructions}


\subsection{Consistently compressed physical potential}

For the particle in the box,
\begin{equation}
    \mathbf H_h=\mathbf T_{\mathrm D,h},
\end{equation}
where $\mathbf T_{\mathrm D,h}$ is the discrete Dirichlet kinetic operator.
Applying the same spectral compression to both terms gives
\begin{equation}
    \mathbf H_h^{(P_M)}
    -
    \mathbf T_{\mathrm D,h}^{(P_M)}
    =
    \mathbf 0.
\end{equation}
This is the consistently reduced physical potential. It vanishes because the
confinement is already included in the Dirichlet realization of the kinetic
operator.

\subsection{Projected Hamiltonian minus unprojected kinetic operator}

Consider instead
\begin{equation}
    \widetilde{\mathbf R}_h
    =
    \mathbf H_h^{(P_M)}
    -
    \mathbf T_{\mathrm D,h}.
\end{equation}
Because $\mathbf P_M$ is a spectral projector of
$\mathbf T_{\mathrm D,h}$, it commutes with that matrix. With
\begin{equation}
    \mathbf Q_M=\mathbf I-\mathbf P_M,
\end{equation}
one obtains
\begin{equation}
    \widetilde{\mathbf R}_h
    =
    -\mathbf Q_M
    \mathbf T_{\mathrm D,h}
    \mathbf Q_M.
\end{equation}
The nonzero residual is therefore the negative discarded kinetic sector. It
results from applying the retention map to the Hamiltonian but not to the
kinetic operator. It is neither the consistently projected potential nor a
boundary operator.

\subsection{Difference between discrete boundary realizations}

Let
\begin{equation}
    \mathbf H_{\mathrm{box},h},
    \mathbf T_{\mathrm{ref},h}
    \in
    \mathcal L(\mathcal V_h)
\end{equation}
be two discrete operators acting on an explicitly identified numerical state
space $\mathcal V_h$. Suppose that $\mathbf H_{\mathrm{box},h}$ carries the
box boundary closure and $\mathbf T_{\mathrm{ref},h}$ carries a different
declared closure. Then
\begin{equation}
    \mathbf B_h
    =
    \mathbf H_{\mathrm{box},h}
    -
    \mathbf T_{\mathrm{ref},h}
\end{equation}
is a well-defined finite-dimensional operator. Depending on the
discretization, $\mathbf B_h$ may be expressed through altered boundary rows,
modified stencils, eliminated degrees of freedom, or nonlocal couplings. In
this construction the residual encodes boundary structure relative to the
selected reference realization.

The identification of the two numerical state spaces is part of the
construction. The result should not be described as a
representation-independent continuum potential.

\section{Residual identifiability}

The three cases demonstrate that the symbol
\begin{equation}
    \mathbf H_{\mathrm{eff}}-\mathbf T
\end{equation}
does not have a unique physical interpretation. Its meaning depends on which
operator was reduced, which reference was selected, and whether both terms
were passed through the same map.

More generally, define
\begin{equation}
    \mathbf R_{T_0}
    =
    \mathbf H_{\mathrm{red}}-\mathbf T_0.
\end{equation}
For any Hermitian matrix $\mathbf K$ on the retained space,
\begin{equation}
    \mathbf H_{\mathrm{red}}
    =
    (\mathbf T_0+\mathbf K)
    +
    (\mathbf R_{T_0}-\mathbf K).
\end{equation}
The reduced Hamiltonian alone therefore does not identify a unique
kinetic--potential decomposition. A potential interpretation requires prior
restrictions on the kinetic reference and on the admissible potential class.

Let
\begin{equation}
    \mathcal M_V
    =
    \left\{
        \mathbf V(\boldsymbol\theta)
        \;\middle|\;
        \boldsymbol\theta\in\Theta
    \right\}
\end{equation}
be a prescribed model class. The best admissible potential is
\begin{equation}
    \mathbf V_{\mathcal M}^{\star}
    =
    \operatorname*{arg\,min}_{\mathbf V\in\mathcal M_V}
    \left\lVert
        \mathbf R_{T_0}-\mathbf V
    \right\rVert_{\mathrm F}^2,
\end{equation}
with unexplained residual
\begin{equation}
    \mathbf E_{\mathcal M}
    =
    \mathbf R_{T_0}
    -
    \mathbf V_{\mathcal M}^{\star}.
\end{equation}
If arbitrary Hermitian matrices are admitted, the minimization is trivial. The
inverse problem becomes scientifically informative only when $\mathcal M_V$ is
restricted by locality, range, symmetry, differential form, parameter sharing,
or transferability.

The particle in a box therefore supplies a controlled demonstration of the
central methodological point: operator subtraction is straightforward after a
common numerical state space has been established, but the physical category
assigned to the result is determined by the reference operator and the
admissible model class.
